% chapters/implementation/technical_details.tex
% Sección 5.3: Aspectos técnicos relevantes (~2-3 páginas)
% ============================================================================
%
% GUÍA: Aspectos transversales de implementación que no pertenecen a un
%       único subsistema. Centrarse en lo que aporta valor técnico y
%       demuestra competencia de ingeniería.
%
%   \subsection{Configuración del entorno con Docker Compose}
%   - Descripción del docker-compose.yml.
%   - Servicios definidos y sus dependencias.
%   - Variables de entorno y fichero .env.
%   - Referencia a \ref{rnf:infra:compose} y \ref{rnf:infra:services}.
%   - Incluir fragmento del docker-compose.yml si es representativo:
%     \Code[cod:compose]{Docker Compose}
%          {Extracto del fichero \texttt{docker-compose.yml}}
%          {codes/docker_compose_excerpt.yml}{1}{30}{1}{yaml}
%   TODO: Copiar un extracto relevante a codes/docker_compose_excerpt.yml.
%
%   \subsection{Configuración de Django y DRF}
%   - Estructura del proyecto Django (apps, settings, urls).
%   - Settings por entorno (dev/prod) mediante variables de entorno.
%   - Configuración de DRF: paginación, autenticación, permisos, parsers,
%     renderers, excepciones personalizadas.
%   - Configuración de CORS.
%
%   \subsection{Integración con MinIO}
%   - Cliente MinIO / boto3 para Python.
%   - Buckets: estructura y política de nombres.
%   - Streaming de archivos grandes (\ref{rnf:perf:files}).
%   - URLs prefirmadas (presigned URLs) para acceso temporal.
%
%   \subsection{Tareas asíncronas con Celery}
%   - Configuración de Celery con Redis como broker.
%   - Definición de tareas (decorators, retries, timeouts).
%   - Celery Beat para tareas periódicas (retención, limpieza).
%   - Workers stateless e idempotentes (\ref{rnf:scale:stateless}).
%
%   \subsection{Seguridad}
%   - Hash de contraseñas (Argon2/bcrypt) — \ref{rnf:sec:passwords}.
%   - Rate limiting (django-ratelimit o throttling de DRF).
%   - Validación y sanitización de entradas — \ref{rnf:sec:input}.
%   - Configuración HTTPS y cookies seguras — \ref{rnf:sec:https}.
%
%   \subsection{Auditoría}
%   - Modelo AuditLog y middleware/señales para registro automático.
%   - Qué eventos se registran.
%   - Referencia a \ref{rnf:gdpr:audit}.
% ============================================================================

% TODO: Redactar los aspectos técnicos relevantes.
